\section{Комбинаторы парсеров}
\label{sec:Combinators}
\tikzsetfigurename{ClassicalParsing_Combinators_}

До сих пор мы рассматривали алгоритмы синтаксического анализа как отдельные процедуры или как интерпретаторы таблиц, построенных по грамматике. Комбинаторы парсеров предлагают другой взгляд: парсер является обычным значением в языке программирования, а более сложные парсеры строятся из простых с помощью функций-комбинаторов~\sidecite{hutton1996monadic}. В результате описание грамматики становится частью программы, и для него доступны обычные средства языка: функции, параметры, модули, обобщённые типы и статическая проверка типов.

На интуитивном уровне парсер для входного алфавита $\Sigma$ и результата типа $A$ можно рассматривать как функцию, которая получает входную строку и возвращает множество пар $(a,w')$, где $a \in A$~--- значение, построенное при разборе, а $w' \in \Sigma^*$~--- ещё не обработанный суффикс входа. Множество результатов нужно для поддержки неоднозначности: один и тот же префикс может быть разобран несколькими способами. В детерминированных практических библиотеках вместо множества часто возвращается один результат или сообщение об ошибке, но математически удобно держать в уме именно недетерминированную модель.

Такой подход занимает промежуточное положение между рукописными парсерами и генераторами парсеров. Как и рукописный парсер, программа на комбинаторах легко встраивается в остальной код и может сразу выполнять семантические действия. Как и грамматика для генератора, композиция комбинаторов обычно близка к форме правил языка и потому остаётся декларативной. Цена этого удобства~--- необходимость аккуратно решать вопросы производительности, левой рекурсии и качества сообщений об ошибках.

\subsection{Базовые парсеры и комбинаторы}
\label{subsec:ParserCombinatorsBasic}

Комбинаторная библиотека обычно предоставляет небольшой набор примитивных парсеров и операций, из которых строятся остальные конструкции. Обозначения различаются от языка к языку, но основные идеи стабильны.

\begin{table}[h]
    \centering
    \caption{Типичные примитивы и комбинаторы парсеров}
    \label{tbl:ParserCombinatorsBasic}
    \begin{tabular}{l|p{0.62\textwidth}}
        \hline
        Обозначение & Смысл \\
        \hline
        $\mathit{fail}$ & всегда завершается неудачей \\
        $\mathit{pure}(a)$ & не читает вход и возвращает значение $a$ \\
        $\mathit{token}(t)$ & читает один токен $t$ \\
        $\mathit{satisfy}(p)$ & читает один токен, удовлетворяющий предикату $p$ \\
        $p \mathbin{\sim} q$ & последовательность: сначала $p$, затем $q$ \\
        $p \mid q$ & альтернатива: разобрать $p$ или $q$ \\
        $p?$ & опциональный парсер: $p$ или пустой разбор \\
        $p^*$ & повторение ноль или более раз \\
        $p^+$ & повторение один или более раз \\
        $\mathit{map}(p,f)$ & разобрать $p$ и применить функцию $f$ к результату \\
        $\mathit{bind}(p,f)$ & разобрать $p$, затем выбрать следующий парсер по полученному результату \\
        \hline
    \end{tabular}
\end{table}

Последовательность и альтернатива соответствуют двум основным операциям над языками: конкатенации и объединению. Если $L(p)$ и $L(q)$~--- языки, распознаваемые парсерами $p$ и $q$, то $p \mathbin{\sim} q$ распознаёт $L(p)L(q)$, а $p \mid q$ распознаёт $L(p) \cup L(q)$. Повторение $p^*$ соответствует замыканию Клини. Отличие от чистой теории языков состоит в том, что парсер обычно возвращает не только факт принадлежности, но и значение: дерево разбора, абстрактное синтаксическое дерево, число, токен или произвольный объект программы.

\begin{example}[Простой парсер выражения]
    Пусть требуется разобрать цепочки вида $a^n b^n$, $n \geq 0$, и вернуть число $n$. Контекстно-свободная грамматика для этого языка имеет вид
    \[
        S \to a S b \mid \varepsilon.
    \]
    В стиле комбинаторов это можно записать следующим образом:
    \begin{verbatim}
lazy val s: Parser[Int] =
  (token("a") ~ s ~ token("b")).map { case _ ~ n ~ _ => n + 1 } |
  pure(0)
    \end{verbatim}
    Здесь рекурсивное значение \texttt{s} соответствует нетерминалу $S$, комбинатор \texttt{\textasciitilde} задаёт последовательность, \texttt{|}~--- альтернативу, а \texttt{map} является семантическим действием, которое преобразует дерево разбора в число. Ключевое слово \texttt{lazy} подчёркивает, что рекурсивный парсер не должен немедленно вычисляться при определении.
\end{example}

В реальных библиотеках поверх базовых операций строятся более удобные комбинаторы: разбор списков с разделителями, скобочных конструкций, ключевых слов, приоритетов операторов, пробелов и комментариев. Именно эта возможность постепенно наращивать предметно-специфичный язык описания синтаксиса является одной из причин популярности комбинаторного подхода.

\subsection{Семантические действия и типы}
\label{subsec:ParserCombinatorsSemantics}

Семантическое действие~--- это функция, применяемая к результату разбора некоторой конструкции. В генераторах парсеров такие действия часто записываются внутри правил грамматики. В комбинаторном подходе они являются обычными функциями языка программирования и комбинируются с парсерами с помощью операций вроде \texttt{map} или \texttt{bind}.

Например, парсер числа может сначала разобрать непустую последовательность цифр, а затем применить к ней функцию преобразования строки в целое число. Парсер арифметического выражения может сразу строить абстрактное синтаксическое дерево. Парсер конфигурационного файла может возвращать готовую структуру данных приложения. При этом тип результата парсера фиксируется в типе самого значения парсера: \texttt{Parser[Int]}, \texttt{Parser[Expr]}, \texttt{Parser[Config]} и т.д.

Статическая типизация помогает обнаруживать ошибки в семантических действиях до запуска программы. Если парсер левой части последовательности возвращает число, а действие пытается использовать его как строку, компилятор языка обнаружит несоответствие типов. Это важное отличие от строковых спецификаций грамматик, в которых ошибки в действиях часто проявляются только при генерации или выполнении парсера.

\subsection{Аппликативные, монадические и селективные комбинаторы}
\label{subsec:ParserCombinatorsKinds}

Комбинаторы парсеров различаются по тому, насколько сильно следующий шаг разбора может зависеть от уже полученного результата. В аппликативном стиле структура разбора известна заранее: парсер строится как статическая композиция независимых частей. Это удобно для анализа грамматики, оптимизации, построения сообщений об ошибках и предварительной проверки свойств.

Монадический стиль использует операцию \texttt{bind}: следующий парсер выбирается функцией от результата предыдущего. Это делает подход выразительнее. Например, можно сначала разобрать число $n$, а затем разобрать ровно $n$ следующих элементов. Но такая динамическая зависимость затрудняет статический анализ: до выполнения первого парсера не всегда известно, какой парсер будет применён дальше.

Селективные комбинаторы занимают промежуточное положение между аппликативными и монадическими. Они позволяют выражать ограниченную форму зависимости от результата без полного перехода к монадической динамике. Современная библиотека Parsley использует эту идею вместе со staging: описание парсера представляется как дерево, частично специализируется и компилируется в эффективный код~\sidecite{10.1145/3409002}. Такой подход показывает, что комбинаторы парсеров не обязаны быть существенно медленнее рукописных или сгенерированных парсеров, если библиотека сохраняет достаточно статической информации о структуре разбора.

\subsection{Левая рекурсия, неоднозначность и мемоизация}
\label{subsec:ParserCombinatorsLeftRecursion}

Наивная реализация комбинаторов парсеров плохо работает с левой рекурсией. Рассмотрим правило
\[
    E \to E + T \mid T.
\]
Если записать его напрямую как рекурсивное определение парсера $E$, то при попытке разобрать $E$ парсер сначала снова вызовет $E$ на той же позиции входа, затем ещё раз и так далее. В обычном рекурсивном спуске это приводит к немедленной бесконечной рекурсии.

Есть несколько способов справиться с этой проблемой. Часто грамматику переписывают, устраняя левую рекурсию вручную. Для арифметических выражений также используют специальные комбинаторы для операторных приоритетов. Однако эти решения не всегда удобны: они меняют структуру грамматики и могут усложнять семантические действия.

Более общий путь~--- мемоизация. Если результат разбора нетерминала в данной позиции входа сохраняется, то повторный запрос к той же паре <<нетерминал, позиция>> не запускает вычисление заново. В packrat-парсерах мемоизация используется для гарантии линейного времени для грамматик с упорядоченным выбором~\sidecite{Johnson1995memoization}. Для произвольных КС-грамматик, неоднозначности и левой рекурсии нужны более сложные механизмы, близкие к обобщённому LL-анализу.

Библиотека Meerkat является примером такого решения~\sidecite{izmaylova2016practical}. Она сочетает идеи комбинаторов парсеров с GLL-подходом~\sidecite{Scott:2010:GP:1860132.1860320}: использует мемоизацию, передачу продолжений и компактное представление результатов. Благодаря этому Meerkat поддерживает произвольные КС-грамматики, включая леворекурсивные и неоднозначные. Это важно не только для классического синтаксического анализа строк, но и для дальнейшего применения комбинаторов в задачах анализа графов.

Неоднозначность требует отдельного представления результата. Если вход имеет несколько деревьев разбора, парсер должен либо вернуть все результаты, либо представить их компактно. В обобщённых алгоритмах для этого используется SPPF, рассмотренное в разделе~\ref{sec:SPPF}. Комбинаторная библиотека может скрывать эту структуру за интерфейсом ленивого перечисления результатов, но внутри ей всё равно необходимо избегать экспоненциального дублирования одинаковых поддеревьев.

\subsection{Преимущества и ограничения}
\label{subsec:ParserCombinatorsTradeoffs}

Главное преимущество комбинаторов парсеров~--- композиционность. Грамматику можно строить из маленьких независимых фрагментов, переиспользовать общие части и параметризовать типовые шаблоны. Кроме того, спецификация синтаксиса находится в том же языке, что и остальная программа, поэтому не требуется отдельный генератор, отдельный синтаксис описания грамматик и отдельный этап сборки.

Второе преимущество~--- естественная интеграция семантики. Парсер сразу возвращает значение нужного типа, а семантические действия являются обычными функциями. Это удобно при построении абстрактных синтаксических деревьев, интерпретаторов, статических анализаторов и предметно-ориентированных языков.

Ограничения также существенны. Наивные комбинаторы могут быть медленными из-за возвратов и повторных вычислений. Хорошие сообщения об ошибках требуют специальной поддержки: простая альтернатива $p \mid q$ может скрыть, какая ветвь действительно продвинулась дальше и чего ожидал парсер в точке ошибки. Левая рекурсия и неоднозначность требуют мемоизации, GLL-подобных механизмов или ограничений на класс грамматик. Наконец, слишком выразительный монадический интерфейс может мешать статической оптимизации.

Таким образом, комбинаторы парсеров являются не заменой всем остальным методам синтаксического анализа, а отдельным инженерным компромиссом. Они особенно полезны, когда важны встраивание в язык общего назначения, композиция, типобезопасность и гибкие семантические действия. Если же требуется максимальная производительность для фиксированной грамматики, генераторы парсеров или специализированные реализации LL/LR/GLL могут быть предпочтительнее.

В дальнейшем мы вернёмся к этой технике в разделе~\ref{sec:CFPQ_Combinators}, где комбинаторы применяются уже не к линейному входу, а к графу. Там примитивный парсер токена заменяется переходом по ребру графа, а результатом разбора становится множество путей или его компактное представление.
